% =============================================================================
% Title Block (spans both columns via \twocolumn[...])
% =============================================================================
\twocolumn[{
\begin{flushleft}
{\small
Department of Software Engineering\\
Faculty of Sciences and Techniques Tangier\\
Abdelmalek Essadi University, Morocco\\
22/03/2026\par

\vspace{2.5em}

{\color{titleblue}\Huge\bfseries Détection des fuites d'eau dans les réseaux\\[0.2em]
urbains de distribution par apprentissage profond\\[0.2em]
avec une architecture CNN-BiLSTM et auto-attention\\[0.2em]
\par}

\vspace{1.5em}

{\large\bfseries Encadré par : Pr. Mostafa EZZIYANI\\[0.8em]
Réalisé par : ZARQI Ezzoubair, El Gharib Mahmoud\par}

}
\end{flushleft}

\vspace{1.5em}

% Abstract
\begin{leftbar}
\textbf{\textcolor{headingblue}{RÉSUMÉ}} Cet article présente un cadre complet
d'apprentissage profond pour la détection automatique des fuites dans les réseaux
urbains de distribution d'eau à partir de séries temporelles multivariées issues
de capteurs. L'approche proposée intègre une ingénierie de caractéristiques guidée
par le domaine, une modélisation séquentielle par fenêtre glissante, et une
architecture hybride combinant des réseaux de neurones convolutionnels (CNN),
une mémoire longue à court terme bidirectionnelle (BiLSTM) et un mécanisme
d'auto-attention.

Le composant CNN capture les motifs temporels locaux, le BiLSTM modélise les
dépendances de long terme, et le mécanisme d'attention met en évidence les pas
de temps les plus informatifs dans chaque séquence. En outre, une stratégie de
séparation temporelle sans fuite de données et une calibration du seuil de
décision sont utilisées afin d'assurer une évaluation robuste et réaliste en
présence de déséquilibre de classes.

Les résultats expérimentaux montrent que la méthode proposée obtient de très
bonnes performances, avec une AUC-PR de $0.9583$ et un score F1 de $0.9318$,
tout en maintenant un rappel élevé. De plus, les comparaisons de benchmark face
à des modèles classiques de machine learning et à des architectures alternatives
d'apprentissage profond confirment l'efficacité du modèle hybride proposé.

Globalement, ce cadre propose une solution pratique et évolutive pour les
systèmes réels de surveillance intelligente de l'eau, contribuant à améliorer la
fiabilité des infrastructures et à réduire les pertes d'eau.
\end{leftbar}

\vspace{1em}

% Keywords
\begin{leftbar}
\textbf{\textcolor{headingblue}{MOTS-CLÉS}} Détection des fuites d'eau,
Apprentissage profond, CNN-BiLSTM, Auto-attention,
Séries temporelles multivariées, Fenêtre glissante,
Classification déséquilibrée, Calibration du seuil,
Comparaison de benchmark, Réseaux d'eau intelligents;
\end{leftbar}

\vspace{2.5em}
}]